% !TEX root = ../main.tex
Neutrinos are elementary particles of the Standard Model (SM). They are electrically neutral leptons that come 
in three flavours, corresponding to the electron, muon and tau charged leptons. Neutrinos have been observed to oscillate or change between these flavours as they propagate, a fact that implies that neutrinos, contrary to the SM prediction, are massive. This chapter begins with a brief history of the neutrino, from its postulate to the discovery of the three flavours and the observation of oscillations. It then places the neutrino in the SM, developing the theory of flavour mixing and oscillations, and reviews the experimental landscape, including three remaining open questions, namely, the mass ordering, the octant of $\theta_{23}$, and the value of $\dcp$.

\section{A Brief History of the Neutrino}
\label{sec:history}

\subsection{The Neutrino Postulate}

By the early 20th century, decay products of atomic nuclei were understood to produce discrete energy spectra. 
This behaviour was observed in both $\alpha$- and $\gamma$-decays:
\begin{align}
    {}^{A}_{Z}X &\rightarrow {}^{A-4}_{Z-2}Y + {}^{4}_{2}\mathrm{\alpha},
    \label{eq:alpha-decay} \\
    {}^{A}_{Z}X^{*} &\rightarrow {}^{A}_{Z}X + \gamma.
    \label{eq:gamma-decay}
\end{align}
In $\alpha$-decay (Eq.~\ref{eq:alpha-decay}), the fixed energy of the emitted $\alpha$ particle is determined by the two-body kinematics of the decay process.
Similarly, the discrete energy spectrum of $\gamma$-decay (Eq.~\ref{eq:gamma-decay}) arises from the de-excitation of the nucleus to a lower energy state, where the difference between the two energy levels corresponds to the emitted photon energy.
This understanding therefore left James Chadwick puzzled when he discovered, in 1914, that the energy spectrum of the emitted electron in $\beta$-decay, then assumed to be the two-body process
\begin{equation}
    {}^{A}_{Z}X \rightarrow {}^{A}_{Z+1}Y + e^{-},
    \label{eq:beta-decay}
\end{equation}
was continuous~\cite{Chadwick1914,Zuber2020}.
The implications of this observation were only addressed in 1930 by Wolfgang Pauli, who proposed that there must be another neutral particle in the final state, so that $\beta$-decay proceeds as
\begin{equation}
    {}^{A}_{Z}X \rightarrow {}^{A}_{Z+1}Y + e^{-} + \bar{\nu}_e.
    \label{eq:beta-decay-threebody}
\end{equation}
In addition to resolving the apparent violation of energy conservation, Pauli's postulate also provided a solution to the fact that the decay could not just emit one spin-$\frac{1}{2}$ particle (a fermion), but must emit another.
Pauli provisionally named this neutral fermion the neutron; it was later renamed the neutrino by Fermi. 

\subsection{First Direct Neutrino Detection and Discovery}
The first direct detection of the electron antineutrino was made in 1956 by Clyde Cowan and Frederick Reines~\cite{ReinesCowan1953,CowanReines1956}.
They used a water tank and layers of liquid scintillators, placed near a nuclear reactor, which provided a large antineutrino flux.
An incoming antineutrino interacting with a proton produces a neutron and a positron via inverse beta decay:
\begin{equation}
    p + \overline{\nu}_e \rightarrow n + e^{+}.
    \label{eq:inverse-beta-decay}
\end{equation}
Two gamma rays are produced when the positron annihilates with an
electron. The detection of these gammas, followed shortly (a few
$\mu$s) by detection of the neutron capture photon, provided a clear
signal of the neutrino-proton interaction.
Reines was awarded the Nobel Prize in Physics in 1995 for the discovery of the electron antineutrino.

\subsection{$\nu_{\mu}$ and $\nu_{\tau}$ Discovery}
Further to Reines and Cowan's discovery, it was understood that reactor and $\beta$-decay neutrinos were associated with an electron (or a $\beta$ particle)~\cite{CowanReines1956} (Eq.~\ref{eq:inverse-beta-decay}). 
In contrast, the question remained as to whether the neutrino resulting from pion decay, produced alongside a muon, was the same neutrino as that of $\beta$-decay~\cite{Feinberg1958,Pontecorvo1960}. 
To answer this, in 1962 Lederman, Schwartz, Steinberger and collaborators at Brookhaven National Laboratory used the world's first accelerator-based neutrino beam~\cite{Danby1962}. 
Accelerated protons were collided with a target to produce pions, which subsequently
decayed to produce muons and neutrinos. 
The resulting neutrinos were detected, and the interactions were found to produce muons rather than electrons. 
This led to the conclusion that there must be at least two types of
neutrinos: those associated with electrons (reactor and $\beta$-decay)
and those associated with muons (pion decay).
Lederman, Schwartz and Steinberger were awarded the 1988 Nobel Prize in Physics for the discovery of the muon neutrino.

When the tau lepton was discovered in 1975~\cite{Perl1975}, a third associated tau neutrino was expected. 
The Direct Observation of the Nu Tau (DONUT) experiment at Fermilab successfully observed it in 2000~\cite{Kodama2001}. 
This completed the direct observation of the three Standard Model (SM) neutrino flavours.

\subsection{Solar and Atmospheric Anomalies}
\label{sec:history-anomalies}
In parallel to the direct observation of the neutrino, Bruno Pontecorvo theorised in 1957 that neutrinos, like neutral kaons, could oscillate.
This was the first suggestion of neutrino oscillations, but in contrast to today's knowledge, and in analogy with the neutral kaon system ($K^0\leftrightarrow\bar{K}^0$), these oscillations were between neutrinos and antineutrinos~\cite{Pontecorvo1957,Pontecorvo1958}.
In 1962, following the muon neutrino discovery, Maki, Nakagawa and Sakata adapted Pontecorvo's proposal by postulating that the muon and electron neutrinos could mix, the first flavour mixing formulation.
This was further developed by Pontecorvo in 1967, who returned to the problem, formulating two neutrino flavour oscillations~\cite{MakiNakagawaSakata1962,Pontecorvo1967} and suggesting that such oscillations would reduce the solar neutrino flux observed on Earth.
Two and three flavour oscillations will be discussed in more detail in Sec.~\ref{sec:oscillations}.

The following year, Ray Davis and collaborators conducted the Homestake experiment, which used a 615 tonne tank of liquid perchloroethylene deep underground to detect solar neutrinos via the reaction: $\nu_{e} + {}^{37}\mathrm{Cl} \rightarrow {}^{37}\mathrm{Ar} + e^{-}$~\cite{Davis1968,Miramonti2009}.
The measured rate of solar neutrinos was found to be significantly lower than that predicted by the Standard Solar Model (SSM).
This discrepancy was known as the \emph{Solar Neutrino Problem}.
Although neutrino oscillations had already been proposed in the previous year, they were not widely accepted as an explanation for this discrepancy.
The solar neutrino deficit was later confirmed by the Kamiokande, SAGE and GALLEX experiments~\cite{Hirata1990,Abazov1991,Anselmann1992}.

A second discrepancy emerged from the measurement of the atmospheric neutrino flux.
These neutrinos are produced when cosmic rays interact with Earth's
atmosphere, producing charged pions and kaons, which subsequently decay
into muons and muon neutrinos. The muons eventually also decay into
electrons, muon neutrinos and electron neutrinos.
The observed muon neutrino flux was lower than expected.
In 1998, the Super-Kamiokande experiment (SK) reported that this deficit exhibited a dependence on the zenith angle, and therefore the distance travelled by the neutrinos through the Earth.
This path length dependence was the first compelling evidence for neutrino oscillations~\cite{Fukuda1998}.

In 2002, the Sudbury Neutrino Observatory (SNO) showed that the total solar neutrino flux was consistent with the SSM, while showing the same electron neutrino flux as previous experiments claiming a deficit.
The total solar neutrino flux was measured via the neutral-current (NC)
interaction, which does not allow for flavour determination and occurs
for all flavours, while the electron neutrino flux was measured via the
charged-current (CC) interaction by measuring the final-state electron.
This suggested that solar electron neutrinos must have undergone flavour transformation as they travelled through the Sun and to the Earth, hence resolving the Solar Neutrino Problem~\cite{Ahmad2002}.
The full interpretation of the SNO result also requires the matter effect, which will be introduced in Sec.~\ref{sec:matter}.


